% !TEX program = lualatex
% =====================================================================
% andamento-2026-11-09-en.tex -- en-US version of
% `andamento-2026-11-09.tex`, the Solara progress report for
% November 5--9, 2026
%
% PARITY WITH THE PORTUGUESE FILE. This is a translation, not a fork: both
% files carry the same pages, the same figures and the same numbers. A
% change in one belongs in the other. Compile from inside `report/`:
%
%   latexmk -lualatex andamento-2026-11-09-en.tex
%
% DELIBERATE DIVERGENCES FROM THE SOURCE FILE, as in `tarefas-dados-en.tex`:
%
%   `babel` loads `american`, and `\sisetup` overrides the style sheet's
%   decimal comma and thin-space grouping with the en-US period and comma.
%
%   Loanwords italicized in the Portuguese file -- desktop, site, commit,
%   raster, zoom, trackpad, pull request -- are set roman here. The italic
%   marks a word as foreign to the surrounding text; in English it is not.
%
%   "Ponto" in the Portuguese file is the interface's own word here, site.
%
% FORM. One column and one page per figure, as in the source: the figure
% at text width at the top, its explanation below it, in the flow. The
% header is `\cabecalho`, the body of the style's `\abertura` without its
% `\twocolumn`.
%
% NUMBERS come from the demonstration project
% `.temp/demo/solara-demo-2026-11.solara`, from the run of the synthetic
% planning example on Nov 9, and, for Figures 2 and 6, from the author's
% captures of the native macOS window over the Italian store.
% =====================================================================
\documentclass[10pt,a4paper,oneside,onecolumn]{book}

\usepackage{iftex}
\ifPDFTeX\usepackage[utf8]{inputenc}\fi
\usepackage[american]{babel}
\usepackage[a4paper,top=2.0cm,bottom=2.2cm,left=2.5cm,right=2.5cm]{geometry}
\usepackage{estilo-reporte}

\sisetup{output-decimal-marker={.},group-separator={,}}

\raggedbottom
\pagestyle{plain}
\setlength{\parskip}{0.4em}
\setlength{\parindent}{0pt}

\autoria{João Leonardi da Silva Melo}{opensource.leonardi@gmail.com}
\marca{../build/appicon.png}

% The body of `\abertura`, without the `\twocolumn` around it.
\makeatletter
\newcommand{\cabecalho}[2]{%
  \vspace*{-0.4em}%
  \noindent
  \begin{minipage}[b]{\marcalargura}%
    \includegraphics[width=\marcalargura]{\re@marca}%
  \end{minipage}%
  \hspace{\marcarespiro}%
  \begin{minipage}[b]{\dimexpr\textwidth-\marcalargura-\marcarespiro\relax}%
    {\sffamily\bfseries\LARGE\hyphenpenalty=10000\exhyphenpenalty=10000
     \raggedright #1\par}%
    \vspace{0.35em}%
    {\sffamily\fontsize{9}{11}\selectfont\color{tinta2}\raggedright #2\par}%
    \vspace{0.25em}%
    \re@assinatura
  \end{minipage}%
  \vspace{0.7em}\par
  \noindent{\color{destaque}\rule{\textwidth}{0.6pt}}%
  \vspace{1.0em}\par}
\makeatother

% A figure at text width, with its caption, for the top of a page.
\newcommand{\figura}[2]{%
  \noindent\includegraphics[width=\textwidth]{figuras/andamento-2026-11/#1}\par
  \captionof{figure}{#2}\vspace{0.6em}}

% A new page that opens with the figure; the text that follows comes below it.
\newcommand{\paginafigura}[2]{\clearpage\figura{#1}{#2}}

\begin{document}

\cabecalho{Summary of the changes to Solara's interface and system planning}
            {Progress report}%
  \figura{03-map-rio-grande-do-norte.jpg}
         {Demonstration project in Rio Grande do Norte, Brazil, with its two
          sites and two areas over the 220~kV lines, the substations and the
          solar and wind plants of the infrastructure GeoPackage.}

\section*{An interface organized around the map}

The map now fills the whole window, and every other function is arranged
around it (Figure~1). In the previous version, the functions were spread
over 14 editors in 11 workspaces, a function whose editor was not open was
invisible, and the editors covered the map while the user worked. In the
current version, a rail on the left leads to the screens for the map
(Explore), sites and areas (Sites and areas), the comparison of results
(Compare), system planning (Planning) and the run graph (Run graph). The
right edge of the map holds the products, results and layers panels,
together with the location and measuring tools, and only one panel is open
at a time.

The demonstration project brings the application's six calculation
products together at two sites and two areas in Rio Grande do Norte. A
site is a point of analysis; the project's sites are São Miguel do
Gostoso, on the coast, and Assú, inland. The areas cover Natal, on the
coast, and the Seridó, inland, and the project holds 11 results.

The layers of a GeoPackage file, the geographic data format standardized
by the OGC (Open Geospatial Consortium), enter the project as a reference,
with the active layer, the visible layers and the color of each one. The
demonstration project uses the file \texttt{infra\_brasile.gpkg}, with
substations, 220~kV and 550~kV lines, and solar and wind plants. The
interface background also changed from a warm gray, which read as yellow
next to the map, to a neutral gray.

\paginafigura{italia-geopackage.jpg}
  {The application's start screen over the infrastructure GeoPackage for
   Italy, in the native macOS window, with the substations aggregated into
   counts.}

\section*{Infrastructure layers for Italy}

The layers of a GeoPackage do not depend on the country. Figure~2 shows
the file \texttt{infra\_italia.gpkg} over the satellite imagery of the
basemap. With more than \num{5000} substations in view, the map draws one
circle per cell, with its area proportional to the number of substations,
and the legend gives the scale for 75, 300 and \num{1203} substations;
when the map is zoomed in, each substation is drawn.

The capture was taken in the application's native macOS window and also
shows the start screen, with the new opening photograph, the start
commands and the recent projects.

\paginafigura{05-result-solar-resource.jpg}
  {Solar resource reading at the Assú site.}

\section*{Solar resource at Assú}

Each result is presented as a reading, that is, a page with a title, up to
four indicators and a main panel whose title states the finding. The solar
resource reading combines the 30-year climatology and the 10-year hourly
series of NASA POWER at the site's radiation cell, and computes the yield
of a fixed array with pvlib.

At Assú, global irradiation was \num{2207}~kWh/m² per year, and the
specific yield was \num{1665}~kWh/kWp for a performance ratio (PR) of 0.80.
The optimal tilt was \num{8.5}°, with a gain of \num{0.8}~\% over the
horizontal plane, and global irradiation varied by \num{1.45}~\% between
years. The main panel shows the daily means of each month for global,
direct and diffuse irradiation; November receives 37~\% more global
irradiation than June.

At the project's other site, São Miguel do Gostoso, global irradiation was
\num{2117}~kWh/m² per year, and the specific yield was \num{1609}~kWh/kWp.

\paginafigura{08-natal-solar-terrain.jpg}
  {Solar terrain reading for the Natal area. On the map, the annual
   irradiation of each cell; in the reading, the area by irradiation band.}

\section*{Solar terrain at Natal}

The solar terrain reading computes the annual irradiation on the plane of
the array in each 30~m cell of the Copernicus GLO-30 digital elevation
model, including shading by the horizon. In the Natal area, the mean was
\num{1991}~kWh/m² per year, and the cell with the lowest irradiation
received \num{1664}~kWh/m², 84~\% of the mean. Mean shading by the horizon
was \num{0.37}~\%, with a maximum of \num{16.3}~\%, and the mean slope was
\num{2.3}°.

The histogram shows the area by irradiation band, and the color scale
under the axis is the same as the map's. Most of the area lies close to
the mean, and \num{2.0}~\% of it lies below \num{1934}~kWh/m².

In the Seridó area, the mean was \num{1976}~kWh/m² per year, the minimum
\num{1501}~kWh/m², and the mean slope \num{5.2}°.

\paginafigura{17-planning-build.jpg}
  {Build view of system planning, with the network of the synthetic
   example drawn as a single-line diagram and the list of components.}

\section*{Building a system planning study}

System planning solves economic dispatch and capacity expansion for a
user-defined network with PyPSA, a library for modeling power systems, and
the HiGHS linear solver. The previous version presented the network as five
spreadsheets, the hourly profiles as comma-separated numbers and the result
as a column of six figures and four tables. The current version splits the
work into two views, Build to assemble the study and Results to read the
run.

In the Build view, the study bar holds the mode, the currency, the CO₂ cap,
the horizon and the sources. The network is drawn as a single-line diagram,
with each bus as a bar, its generators above it and its loads and
batteries below. Selecting a component opens it in an inspector, which
shows the availability or demand profile as bars and accepts direct
editing, pasting from a spreadsheet and CSV import.

Before optimization, the draft is checked against the same rules as the
calculation engine, and the Optimize button stays disabled while any fault
remains. The diagram can be panned and zoomed, like the map.

\paginafigura{planning-build-foggia.jpg}
  {Build view of a study with 12 buses in the Foggia region, Italy, in the
   native macOS window.}

\section*{A planning study with a 12-bus network}

The Foggia--Troia study, an expansion with solar generation and batteries,
has 12 buses at 380~kV, 15 generators, 5 loads, 12 lines and 2 batteries.
The horizon starts on June 15, 2026, and has 168 intervals of 1~h. The
existing generators are gas plants, such as the Termoli plant at 770~MW,
wind farms aggregated by substation, and the Troia solar park at 103~MW.
The expansion candidates are two solar plants, of up to 800~MW and up to
400~MW, and two batteries of up to 200~MW with 4~h of storage.

At this size, the diagram arranges the buses in rows and routes the lines
between non-adjacent buses through lanes under each row. The study bar
shows that the draft passed the calculation engine's checks (Ready to
optimize). The component names carry, in brackets, the source identifier
of each object, such as those from OpenStreetMap.

\paginafigura{19-planning-results-dispatch.jpg}
  {Results view of system planning, with the hourly dispatch by
   technology, battery charging below zero and demand as a step line.}

\section*{Reading a system planning run}

In the Results view, four indicators precede a panel with the Dispatch,
Capacity, Lines and Storage tabs. In the synthetic example of two buses
and 24 hours, the horizon cost was \num{71255}~EUR, of which
\num{44453}~EUR was operation and \num{26803}~EUR added capacity. The cost
per MWh served was \num{59.1}~EUR, and direct emissions totaled 177~t of
CO₂.

The optimizer added \num{85.3}~MW of solar generation, 80~MW of battery
and \num{71.7}~MVA to the North–South line, which ran at its thermal limit
for 3~h. Solar generation supplied 64~\% of the energy generated. The
battery charged between 08:00 and 17:00 and discharged 70~\% of its energy
between 17:00 and 21:00, the interval in which demand rises from 45~MW to
70~MW.

\begin{objecao}
The example is fictitious and does not represent a real network. The study
is kept in the application's local storage, outside the \texttt{.solara}
file, and runs last only for the session.
\end{objecao}

\clearpage

\section*{Fix to the connection with the grid database}

Since commit \texttt{24f2bf3}, of October 29, the application's quick
queries have been answered by Python processes kept open between requests.
In these processes, standard output is redirected, and the answer to each
action has to go through the protocol's reply function. The power grid
actions and the map's wind field wrote their answer directly to standard
output, where it was lost. As a result, connecting to any grid database,
called Grid store in the interface, failed with the message ``the action
produced no result'', including the Italian database \texttt{solara\_it}.
Both actions now answer through the protocol, and a new test rejects any
action that writes to standard output outside it.

\end{document}
